\section{Identification and projection bounds}

The causal full-data extension used for interpretation is the following.

The statistical construction uses observable product coordinates to recover the homogeneous conditional log-odds coefficient. Its causal interpretation follows from the potential-outcome coupling in \cref{obj:def:causal-extension}. An admissible coupling for an observed law \(P\) is a Markov kernel \(\Gamma\) from \([0,1]\) to \(\{0,1\}^2\) whose first and second margins, at every covariate value, are the observed control-arm and treated-arm Bernoulli outcome laws. This specifies the conditional joint distribution of the two potential outcomes. The full-data construction then draws the covariate from the uniform law \leanref{sym:\lambda}{\ensuremath{\lambda}}, draws the potential-outcome pair from this coupling, and draws treatment using the observed conditional treatment probability. The following statement gives the observable coordinate identity and the properties of this construction.

\begin{lemmav}{lem:observable-odds}[Observable odds and causal interpretation]
Let \(0<\beta<1/4\) and \(\beta<\alpha\leq1\), and let \(P\in\mathcal M(\alpha,\beta)\), the model class in \cref{obj:def:model}. Write \(\lambda\) for the uniform probability law on \([0,1]\), and define the observable coordinates
\[
C(P)
=
E_P[AY]
-
\int_{[0,1]} e(P,x)\,E_P[Y\mid X=x]\,\lambda(dx),
\qquad
S(P)
=
\int_{[0,1]}
P(A=1,Y=0\mid X=x)\,
P(A=0,Y=1\mid X=x)\,\lambda(dx),
\]
where \(e(P,x)\) is the conditional treatment probability. With
\[
\theta(P)=\operatorname{logit}\mu_1(P,0)-\nu(P,0),
\qquad s_0=\frac{1}{512},
\]
where \(\mu_1(P,0)\) is the treated-arm risk at zero and \(\nu(P,0)\) is the control-arm prognosis logit at zero, one has
\[
C(P)=\{\exp(\theta(P))-1\}S(P),
\qquad
S(P)\geq s_0,
\qquad
\theta(P)=\log\{1+C(P)/S(P)\}.
\]

For every Markov kernel \(\Gamma\) from \([0,1]\) to \(\{0,1\}^2\) whose first and second margins at every \(x\) are Bernoulli laws with means
\[
P(Y=1\mid A=0,X=x)
\quad\text{and}\quad
P(Y=1\mid A=1,X=x),
\]
respectively, the full-data law \(\widetilde P_{P,\Gamma}\) of \cref{obj:def:causal-extension} has the following properties:
\begin{itemize}
\item \textbf{(Observed law and consistency.)}
The image of \(\widetilde P_{P,\Gamma}\) under
\[
((x,(y^0,y^1)),a)\longmapsto
\bigl(x,a,(1-a)y^0+ay^1\bigr)
\]
is \(P\). For every full-data record,
\[
Y=(1-A)Y^0+AY^1.
\]

\item \textbf{(Construction and conditional exchangeability.)}
The law \(\widetilde P_{P,\Gamma}\) is the successive kernel composition obtained by drawing \(X\) from \(\lambda\), drawing \((Y^0,Y^1)\) from \(\Gamma(P,X)\), and drawing \(A\) from the observed-law treatment kernel with Bernoulli mean \(e(P,X)\). Under its full conditional law at every \(x\in[0,1]\), the pair \((Y^0,Y^1)\) and \(A\) are independent.

\item \textbf{(Conditional causal log-odds.)}
For every \(x\in[0,1]\),
\[
\operatorname{logit}
\left(\int_{\{0,1\}^2} y^1\,\Gamma(P,x)(d(y^0,y^1))\right)
-
\operatorname{logit}
\left(\int_{\{0,1\}^2} y^0\,\Gamma(P,x)(d(y^0,y^1))\right)
=
\theta(P).
\]

\item \textbf{(Full conditional disintegration.)}
Composing \(\lambda\) with the full conditional laws of \(((Y^0,Y^1),A)\) at \(x\), then taking the image under the canonical reassociation
\[
\bigl(x,((y^0,y^1),a)\bigr)
\longmapsto
\bigl((x,(y^0,y^1)),a\bigr),
\]
gives exactly \(\widetilde P_{P,\Gamma}\).
\end{itemize}
\end{lemmav}

The coordinate identity in \cref{obj:lem:observable-odds} identifies the effect through an observable ratio, with a denominator bounded uniformly away from zero over the model. Its causal clauses establish that every admissible coupling produces the observed law and the same conditional causal log-odds coefficient. Conditional exchangeability is a property of the specified full-data construction, while the marginal restrictions ensure that its potential-outcome risks equal the observed arm risks.

To control the coordinate estimates in \cref{obj:def:coordinate-estimates}, retain the rank-\(k\) cosine projection \(\Pi_k\) and ordered-pair statistic from \cref{obj:def:projection-statistic}. The approximation bound below concerns a continuous function \leanref{sym:f}{\ensuremath{f}} at Hölder exponent \leanref{sym:\gamma}{\ensuremath{\gamma}}, with Hölder seminorm \leanref{sym:[\cdot]_\gamma}{\ensuremath{[f]_\gamma}}. For two bounded record marks, let \leanref{sym:w}{\ensuremath{w}} and \leanref{sym:v}{\ensuremath{v}} denote their conditional means. The next statement connects approximation of these means to the bias and variance of the product statistic.

\begin{lemmav}{lem:projection-control}[Projection approximation and moments]
Let \(P\) be an observed law with continuous, strictly interior conditional four-cell probabilities, satisfying \cref{obj:ass:uniform-design}. Let \(k\ge1\) be an integer, and let \(\Pi_k\) denote the rank-\(k\) cosine projection used in \cref{obj:def:projection-statistic}.

For every \(0<\gamma\le1\) and every continuous function \(f:[0,1]\to\mathbb R\),
\[
\|f-\Pi_k f\|_{L^2(\lambda)}
\le 5[f]_\gamma k^{-\gamma},
\qquad
[f]_\gamma
=
\sup_{\substack{x,z\in[0,1]\\x\ne z}}
\frac{|f(x)-f(z)|}{|x-z|^\gamma}.
\]
Here \(\lambda\) is the uniform probability law on \([0,1]\), and the norm and Hölder seminorm are interpreted in the extended nonnegative reals, allowing \([f]_\gamma=+\infty\).

For every pair of measurable record marks \(W,V\) taking values in \([0,1]\), and every integer \(n\ge2\), define their conditional means using the four-cell conditional law under \(P\):
\[
w(x)=E_P[W\mid X=x],
\qquad
v(x)=E_P[V\mid X=x].
\]
For the ordered-pair statistic \(\mathbb U_{n,k}(W,V)\) of \cref{obj:def:projection-statistic}, evaluated on \(n\) independent original records with law \(P\),
\[
E_P\mathbb U_{n,k}(W,V)
=
\langle\Pi_k w,\Pi_k v\rangle_{L^2(\lambda)},
\]
\[
\langle w,v\rangle_{L^2(\lambda)}
-
E_P\mathbb U_{n,k}(W,V)
=
\langle w-\Pi_k w,v-\Pi_k v\rangle_{L^2(\lambda)},
\]
and
\[
\operatorname{Var}_P\{\mathbb U_{n,k}(W,V)\}
\le
\frac4n+\frac{2k}{n(n-1)}.
\]
\end{lemmav}

The product-bias identity in \cref{obj:lem:projection-control} expresses the approximation error as the inner product of two projection residuals. It therefore connects the smoothness of both conditional means to the bias of their estimated product. The approximation inequality covers the Lipschitz endpoint with the same constant, and the variance bound separates a term proportional to \(n^{-1}\) from a term proportional to the projection rank divided by the number of ordered pairs. Under the known uniform design, these bounds apply directly in \(L^2(\lambda)\). Together with the denominator floor and inversion in \cref{obj:lem:observable-odds}, they supply the statistical bounds used to calibrate the coordinate error radii in \cref{obj:def:upper-handle} and establish the interval guarantee in \cref{obj:thm:finite-honest-upper}.
